\documentclass[journal]{IEEEtai}
\usepackage{graphicx,amsmath,amssymb,url}
\newtheorem{PWTheorem}{Theorem}
\newcommand{\riskname}{selective risk}
\def\methodname{Gate}
\providecommand{\theHalgorithm}{\arabic{algorithm}}
\title{A Fixed Confidence Gate for Selective Classification}
\author{Anonymous submission}
\markboth{Confidence Gate}{Anonymous submission}
\begin{document}\maketitle
\begin{abstract}
We examine a fixed confidence gate for selective classification using a synthetic
evaluation fixture. The gate abstains from predictions below a specified confidence
threshold. We report accuracy among accepted predictions alongside coverage on two
synthetic datasets. Conditional accuracy increases on one dataset and decreases on
the other; these observations do not establish an overall accuracy improvement.
A range bound for selective risk requires positive coverage and bounded loss.
Statistical uncertainty, runtime, and distribution shift remain unevaluated.
\end{abstract}
\begin{IEEEImpStatement}
Abstaining from uncertain predictions can change both the set of cases a system
handles and the information available to a decision maker. The synthetic fixture
here illustrates why conditional accuracy should be considered together with
coverage: a higher accuracy among accepted cases does not establish better
performance across all cases. This distinction may help researchers design
evaluations that expose the tradeoff between selective prediction and abstention.
The observations do not support deploying the gate in a particular application.
Only a fixed classifier, a single baseline, and synthetic data are considered;
distribution shift, statistical uncertainty, latency, and consequences for affected
people remain unevaluated. Application-specific evidence and an assessment of
rejected cases would be needed before making deployment decisions.
\end{IEEEImpStatement}
\begin{IEEEkeywords}Selective classification, abstention, evaluation, coverage\end{IEEEkeywords}
\section{Introduction}
Selective classification allows a model to abstain from uncertain predictions.
We study a fixed confidence gate, rather than a learned rejection policy.
This synthetic evaluation asks whether a fixed threshold changes accuracy among
accepted predictions. Coverage is reported alongside conditional accuracy because
the gate changes which predictions are evaluated.
\section{Method}\label{sec:method}
Gate uses a fixed classifier and accepts a prediction when its maximum class
probability is at least 0.8; otherwise, it abstains. Always uses the same classifier
and accepts every prediction. Neither method retrains the model.
For a fixed classifier $h$ and acceptance function $a(x)\in\{0,1\}$, assume
$0\leq\ell(h(x),y)\leq1$ and positive coverage $C=\mathbb{E}[a(X)]>0$.
Selective risk is
\begin{equation}\label{eq:risk}
R_{\mathrm{sel}}=\frac{\mathbb{E}[a(X)\ell(h(X),Y)]}{C}.
\end{equation}
\begin{PWTheorem}\label{thm:bounded}
Under these assumptions, $0\leq R_{\mathrm{sel}}\leq1$.
\end{PWTheorem}
The result is a range bound, not a generalization guarantee or a comparison
between Gate and Always. Appendix~\ref{app:proof} gives the full proof.
\section{Synthetic Evaluation}\label{sec:results}
Alpha and Beta each contain 1000 synthetic samples. Both methods use the same
model and 5 evaluation seeds; results are seed means, with no variance estimates.
\begin{table}[htbp]\centering
\caption{Synthetic fixture values. Always accuracy uses all samples;
Gate accuracy is conditional on acceptance. These values are not research findings.}
\label{tab:results}
\begin{tabular}{lrrr}\hline
Dataset & Always (\%) & Gate (\%) & Coverage (\%)\\\hline
Alpha & 80.0 & 84.0 & 60.0\\
Beta & 76.0 & 75.5 & 58.0\\\hline
\end{tabular}\end{table}
Table~\ref{tab:results} shows a higher conditional accuracy on Alpha and a lower
conditional accuracy on Beta. Different evaluation populations prevent interpreting
these differences as an overall accuracy improvement. No significance tests,
runtime measurements, or distribution-shift evaluations were performed.
\section{Limitations and Conclusion}
A single baseline and synthetic datasets do not establish state-of-the-art
performance. The classifier and data-generation details are
missing and required for reproducibility. The fixed threshold does not
guarantee lower risk, and deployment benefits remain untested.
Formatting uses the official author guidelines~\cite{formatguide}.
\label{pw:main-end}

\bibliographystyle{IEEEtran}\bibliography{references}
\appendix
\section{Proof of the Range Bound}\label{app:proof}
Since $0\leq a(X)\ell(h(X),Y)\leq a(X)$, taking expectations gives
$0\leq\mathbb{E}[a(X)\ell(h(X),Y)]\leq C$. Division by $C>0$ proves
Theorem~\ref{thm:bounded}. This argument applies to the specified distribution;
it does not establish a bound under distribution shift.

\end{document}
